% figures17.tex -- Appendix figures redrawn in TikZ.
%
% The Appendix (book pp.275-282) contains exactly ONE figure: the unnumbered
% and uncaptioned diagram on book p.277 that accompanies Postulate 16. It
% shows two three-ray configurations, unprimed and primed, drawn congruent.
% (Pages 275, 276, 278-281 are text only; 282 is blank. Verified page by page
% against sources/Geometry.pdf pp.289-296.)
%
% GEOMETRY MEASURED off sources/Geometry.pdf p.291 rasterised at 600 dpi.
% Method: threshold the ink, separate the heavy rays from the light arcs by
% binary erosion (rays survive a 5x5 erosion, arcs do not), take the vertex as
% the left end of the horizontal ray, then fit ray directions from the eroded
% ray cores only. Reading the whole ink mask instead biases every direction
% toward the arcs that run tangent to it -- that is what produced the earlier
% 45.0/117.0 figures. Raw and de-skewed (photo tilt removed by subtracting the
% measured tilt of the horizontal ray) results:
%
%     ray          raw unprimed  raw primed   de-skewed mean
%     horizontal       +0.48        +0.67          0 (datum)
%     diagonal         44.38        43.87        43.55
%     upper-left      116.33       116.13       115.66
%
% Drawn at 0 / 43.6 / 115.7 deg. Ray lengths and arc radii as fractions of the
% horizontal ray (mean of the two sub-figures):
%
%     horizontal ray 1.000   arc A 0.446   (drawn 3.30 / 1.47)
%     diagonal ray   0.994   arc C 0.699   (drawn 3.28 / 2.31)
%     upper-left ray 0.939   arc B 0.774   (drawn 3.10 / 2.55)
%
% The three arc radii are DISTINCT and ordered A < C < B, and each label sits
% on its own arc in a gap in that arc. Arc spans read off the source by
% scanning each arc's circle for ink (rays and label glyphs masked out):
%
%     arc A   4..60  and  78..114   gap 60..78, label at  69 deg
%     arc B  47..77  and  87..114   gap 77..87, label at  81.5 deg
%     arc C   4..16  and  30..42    gap 16..30, label at  23 deg
%
% The arrowhead ends below are the source's exactly; the GAP ends are opened a
% few degrees wider (A 52..86, B 71.5..91.5, C 12..34) because our labels are
% set with more padding than the book's tight hand lettering -- at the source's
% own gap widths check_labels.py reports the glyph boxes touching the arcs.
%
% Arrowheads therefore land at 114 (upper-left ray), 47/42 (diagonal ray) and
% 4 (horizontal ray) -- always just short of the ray they point at, and the two
% hits on the horizontal ray (arc A at 1.47, arc C at 2.31) and the two on the
% diagonal (arc B at 2.55, arc C at 2.31) are at different radii, as in the book.
%
% LINE WEIGHT: the book draws the three rays roughly 1.75x the weight of the
% annotation arcs (measured 13.0 px vs 8.5 px of ink at 600 dpi on the same
% photo). The rays keep the house `given' weight so the Appendix matches the
% rest of the book; the arcs are dropped to 0.42pt to restore the contrast.
%
% What the postulate rests on: the diagonal ray is INTERIOR to angle A, so
% angle A = angle B + angle C. See tools/constraints/ch17.py.

% ------------------------------------------------ Appendix figure (p.277)
\newcommand{\FIGAPPONE}{\begin{bkfigure}[0.52\linewidth]
% --- unprimed configuration
\begin{tikzpicture}[scale=0.72,
    mk/.style={fig, line width=0.42pt},          % light annotation arcs
    alab/.style={vlab, outer sep=1.5pt}]         % label sits in its arc's gap
  \coordinate (O) at (0,0);
  \draw[given] (O) -- (0:3.30);        % horizontal ray
  \draw[given] (O) -- (43.6:3.28);     % diagonal ray, interior to angle A
  \draw[given] (O) -- (115.7:3.10);    % upper-left ray
  % angle A : upper-left ray <-> horizontal ray  (the whole angle)
  \draw[mk,->] (86:1.47) arc[start angle=86, end angle=114, radius=1.47];
  \draw[mk,->] (52:1.47) arc[start angle=52, end angle=4,   radius=1.47];
  \node[alab] at (69:1.47) {$A$};
  % angle B : upper-left ray <-> diagonal ray
  \draw[mk,->] (91.5:2.55) arc[start angle=91.5, end angle=114, radius=2.55];
  \draw[mk,->] (71.5:2.55) arc[start angle=71.5, end angle=47,  radius=2.55];
  \node[alab] at (81.5:2.55) {$B$};
  % angle C : diagonal ray <-> horizontal ray
  \draw[mk,->] (34:2.31) arc[start angle=34, end angle=42, radius=2.31];
  \draw[mk,->] (12:2.31) arc[start angle=12, end angle=4,  radius=2.31];
  \node[alab] at (23:2.31) {$C$};
\end{tikzpicture}

\vspace{4mm}

% --- primed configuration (drawn congruent to the unprimed one, as the book
%     draws it; the postulate itself only supposes congruence of two pairs)
\begin{tikzpicture}[scale=0.72,
    mk/.style={fig, line width=0.42pt},
    alab/.style={vlab, outer sep=1.5pt}]
  \coordinate (O) at (0,0);
  \draw[given] (O) -- (0:3.30);
  \draw[given] (O) -- (43.6:3.28);
  \draw[given] (O) -- (115.7:3.10);
  \draw[mk,->] (86:1.47) arc[start angle=86, end angle=114, radius=1.47];
  \draw[mk,->] (52:1.47) arc[start angle=52, end angle=4,   radius=1.47];
  \node[alab] at (69:1.47) {$A'$};
  \draw[mk,->] (91.5:2.55) arc[start angle=91.5, end angle=114, radius=2.55];
  \draw[mk,->] (71.5:2.55) arc[start angle=71.5, end angle=47,  radius=2.55];
  \node[alab] at (81.5:2.55) {$B'$};
  \draw[mk,->] (34:2.31) arc[start angle=34, end angle=42, radius=2.31];
  \draw[mk,->] (12:2.31) arc[start angle=12, end angle=4,  radius=2.31];
  \node[alab] at (23:2.31) {$C'$};
\end{tikzpicture}
\end{bkfigure}}
